\begin{tomtat}
	\begin{minipage}[t]{0.3\textwidth}
		\vspace{0pt}
		\begin{tomtat}
			\begin{center}
				\textbf{ĐỀ 55}
			\end{center}
			
		\end{tomtat}
	\end{minipage}
	\hfill
	\begin{minipage}[t]{0.5\textwidth}
		\vspace{0pt}
		\begin{tomtat}
			\begin{center}
				%\includegraphics[scale=0.08]{images/55}
			\end{center}
		\end{tomtat}
	\end{minipage}
\end{tomtat}
\textbf{I. PHẦN TRẮC NGHIỆM}
\begin{ex}
	Kết quả rút gọn biểu thức $A=2\sqrt{18}-5\sqrt{2}$ là
	\choice{$\sqrt{2}$}{$2\sqrt{2}$}{$-\sqrt{2}$}{$3\sqrt{2}$}
	\loigiai{
		\textbf{Đáp án đúng là A.} \\
		Ta có:
		\begin{align*} A &= 2\sqrt{18}-5\sqrt{2} \\ &= 2\sqrt{9 \cdot 2} - 5\sqrt{2} \\ &= 2 \cdot 3\sqrt{2} - 5\sqrt{2} \\ &= 6\sqrt{2} - 5\sqrt{2} \\ &= (6-5)\sqrt{2} \\ &= \sqrt{2} \end{align*}
	}
\end{ex}

\begin{ex}
	Nghiệm của phương trình $4x+1=5x$ là:
	\choice{2}{3}{4}{1}
	\loigiai{
		\textbf{Đáp án đúng là D.} \\
		Ta có:
		\begin{align*} 4x+1 &= 5x \\ 1 &= 5x - 4x \\ 1 &= x \\ x &= 1 \end{align*}
		Vậy, nghiệm của phương trình là $x=1$.
	}
\end{ex}

\begin{ex}
	Cho hàm số $y=f(x)=(-2m+1)x^2$. Tìm giá trị m để đồ thị đi qua điểm A(-2;4).
	\choice{m=0}{m=1}{m=2}{m=-2}
	\loigiai{
		\textbf{Đáp án đúng là A.} \\
		Để đồ thị hàm số $y=(-2m+1)x^2$ đi qua điểm A(-2;4), tọa độ của điểm A phải thỏa mãn phương trình hàm số:
		$4 = (-2m+1)(-2)^2$
		\begin{align*} 4 &= (-2m+1) \cdot 4 \\ 1 &= -2m+1 \\ 2m &= 1-1 \\ 2m &= 0 \\ m &= 0 \end{align*}
		Vậy, $m=0$.
	}
\end{ex}

\begin{ex}
	Một hộp có 20 quả bóng được đánh số từ 1 đến 20, đồng thời các quả bóng từ 1 đến 5 được sơn màu cam và các quả bóng còn lại được sơn màu xanh; các quả bóng có kích cỡ và khối lượng như nhau. Lấy ngẫu nhiên một quả bóng trong hộp. Không gian mẫu có bao nhiêu phần tử?
	\choice{30}{10}{20}{15}
	\loigiai{
		\textbf{Đáp án đúng là C.} \\
		Khi lấy ngẫu nhiên một quả bóng trong hộp có 20 quả bóng, mỗi quả bóng có khả năng được chọn như nhau.
		Không gian mẫu là tập hợp tất cả các kết quả có thể xảy ra.
		Vì có 20 quả bóng, nên có 20 kết quả có thể xảy ra khi lấy ngẫu nhiên một quả.
		Vậy, không gian mẫu có 20 phần tử.
	}
\end{ex}
\begin{ex}
	Cho tam giác ABC vuông tại A, có AB = 3 cm, AC = 4 cm. $\sin B$ bằng
	\choice{$\dfrac{4}{5}$}{$\dfrac{4}{3}$}{$\dfrac{3}{4}$}{$\dfrac{3}{5}$}
	\loigiai{
		\textbf{Đáp án đúng là A.} \\
		Xét tam giác ABC vuông tại A, theo định lý Pythagoras ta có:
		$BC^2 = AB^2 + AC^2 = 3^2 + 4^2 = 9 + 16 = 25$
		$\Rightarrow BC = \sqrt{25} = 5$ cm.
		Trong tam giác vuông ABC, ta có:
		$\sin B = \dfrac{\text{cạnh đối}}{\text{cạnh huyền}} = \dfrac{AC}{BC} = \dfrac{4}{5}$.
	}
\end{ex}

\begin{ex}
	Tính thể tích của một hình trụ có bán kính đáy 3 cm, chiều cao 8 cm là
	\choice{$24\pi \text{ cm}^3$}{$72\pi \text{ cm}^3$}{$192\pi \text{ cm}^3$}{$12\pi \text{ cm}^3$}
	\loigiai{
		\textbf{Đáp án đúng là B.} \\
		Thể tích của hình trụ được tính theo công thức $V = \pi r^2 h$, trong đó $r$ là bán kính đáy và $h$ là chiều cao.
		Thay $r = 3$ cm và $h = 8$ cm vào công thức, ta có:
		$V = \pi (3)^2 \cdot 8 = \pi \cdot 9 \cdot 8 = 72\pi \text{ cm}^3$.
	}
\end{ex}

\begin{ex}
	Nghiệm của bất phương trình $2x+1 \ge x+2$ là
	\choice{$x<1$}{$x>1$}{$x \ge 1$}{$x \le 1$}
	\loigiai{
		\textbf{Đáp án đúng là C.} \\
		Ta có:
		\begin{align*} 2x+1 &\ge x+2 \\ 2x - x &\ge 2 - 1 \\ x &\ge 1 \end{align*}
		Vậy, nghiệm của bất phương trình là $x \ge 1$.
		(Lưu ý: Trong hình ảnh đề bài, đáp án D là $x \ge 1$, có vẻ như có sự nhầm lẫn giữa C và D trong các lựa chọn của đề gốc. Dựa trên giải, $x \ge 1$ là đúng).
	}
\end{ex}
\begin{ex}
	Sau buổi kiểm tra môn Toán, giáo viên ghi lại số điểm mà học sinh đạt được như sau:
	\begin{center}
		\begin{tabular}{|*{20}{c|}}
			\hline
			5 & 6 & 8 & 7 & 4 & 5 & 7 & 6 & 8 & 9 & 7 & 5 & 6 & 8 & 7 & 8 & 8 & 9 & 8 & 7 \\
			\hline
			10 & 6 & 5 & 6 & 5 & 4 & 5 & 6 & 10 & 7 & 6 & 5 & 7 & 6 & 10 & 8 & 5 & 5 & 6 & 5 \\
			\hline
		\end{tabular}
	\end{center}
	Tần số của giá trị 10 điểm là
	\choice{4}{1}{2}{3}
	\loigiai{
		\textbf{Đáp án đúng là D.} \\
		Để tìm tần số của giá trị 10 điểm, ta đếm số lần xuất hiện của điểm 10 trong bảng số liệu đã cho (đã được sắp xếp lại thành 2 hàng).
		Quan sát bảng số liệu:
		\begin{itemize}
			\item Hàng 1: Không có điểm 10 nào.
			\item Hàng 2: Điểm 10 xuất hiện ở các vị trí cột thứ 1, cột thứ 9, và cột thứ 15. Do đó, có 3 điểm 10.
		\end{itemize}
		Tổng số lần xuất hiện của điểm 10 là $0 + 3 = 3$.
		Vậy, tần số của giá trị 10 điểm là 3.
	}
\end{ex}

\newpage
\textbf{II. PHẦN TỰ LUẬN}
\begin{ex}
	Rút gọn biểu thức: $B = \left(\dfrac{2\sqrt{x}+3}{x-9} - \dfrac{1}{\sqrt{x}+3}\right) : \dfrac{\sqrt{x}+6}{\sqrt{x}+3}$ với $x \ge 0; x \ne 9$.
	\loigiai{
		Với $x \ge 0$ và $x \ne 9$, ta có:
		\begin{align*} B &= \left(\dfrac{2\sqrt{x}+3}{(\sqrt{x}-3)(\sqrt{x}+3)} - \dfrac{1}{\sqrt{x}+3}\right) : \dfrac{\sqrt{x}+6}{\sqrt{x}+3} \\ &= \left(\dfrac{2\sqrt{x}+3}{(\sqrt{x}-3)(\sqrt{x}+3)} - \dfrac{\sqrt{x}-3}{(\sqrt{x}-3)(\sqrt{x}+3)}\right) \cdot \dfrac{\sqrt{x}+3}{\sqrt{x}+6} \\ &= \dfrac{2\sqrt{x}+3 - (\sqrt{x}-3)}{(\sqrt{x}-3)(\sqrt{x}+3)} \cdot \dfrac{\sqrt{x}+3}{\sqrt{x}+6} \\ &= \dfrac{2\sqrt{x}+3 - \sqrt{x}+3}{(\sqrt{x}-3)(\sqrt{x}+3)} \cdot \dfrac{\sqrt{x}+3}{\sqrt{x}+6} \\ &= \dfrac{\sqrt{x}+6}{(\sqrt{x}-3)(\sqrt{x}+3)} \cdot \dfrac{\sqrt{x}+3}{\sqrt{x}+6}\end{align*}
		Vì $x \ge 0$ nên $\sqrt{x} \ge 0$, do đó $\sqrt{x}+6 > 0$.
		Vì $x \ne 9$ nên $\sqrt{x} \ne 3$, do đó $\sqrt{x}+3 > 0$.
		\begin{align*} B &= \dfrac{1}{\sqrt{x}-3} \end{align*}
		Vậy, $B = \dfrac{1}{\sqrt{x}-3}$ với $x \ge 0; x \ne 9$.
	}
\end{ex}

\begin{ex}
	Cho phương trình: $x^2-9x+16=0$ có hai nghiệm dương phân biệt $x_1, x_2$. Không giải phương trình, hãy tính giá trị của biểu thức $T = \dfrac{x_1\sqrt{x_2}+x_2\sqrt{x_1}}{x_1^2+x_2^2}$.
	\loigiai{
		Phương trình $x^2-9x+16=0$ có:
		$\Delta = (-9)^2 - 4 \cdot 1 \cdot 16 = 81 - 64 = 17 > 0$.
		Vì $\Delta > 0$, phương trình có hai nghiệm phân biệt.
		Theo hệ thức Vi-ét, ta có:
		$\begin{cases} x_1+x_2 = -\dfrac{-9}{1} = 9 \\ x_1x_2 = \dfrac{16}{1} = 16 \end{cases}$
		Vì $x_1+x_2 = 9 > 0$ và $x_1x_2 = 16 > 0$, nên phương trình có hai nghiệm dương phân biệt $x_1, x_2$.
		Điều này đảm bảo $\sqrt{x_1}$ và $\sqrt{x_2}$ có nghĩa.
		
		Xét tử số của T:
		$x_1\sqrt{x_2}+x_2\sqrt{x_1} = \sqrt{x_1x_2}(\sqrt{x_1}+\sqrt{x_2})$.
		Ta có $(\sqrt{x_1}+\sqrt{x_2})^2 = x_1+x_2+2\sqrt{x_1x_2}$.
		Vì $x_1, x_2$ dương, $\sqrt{x_1}+\sqrt{x_2} > 0$.
		Nên $\sqrt{x_1}+\sqrt{x_2} = \sqrt{x_1+x_2+2\sqrt{x_1x_2}}$.
		Thay số: $\sqrt{x_1}+\sqrt{x_2} = \sqrt{9+2\sqrt{16}} = \sqrt{9+2 \cdot 4} = \sqrt{9+8} = \sqrt{17}$.
		Vậy, tử số là $\sqrt{16} \cdot \sqrt{17} = 4\sqrt{17}$.
		
		Xét mẫu số của T:
		$x_1^2+x_2^2 = (x_1+x_2)^2 - 2x_1x_2$.
		Thay số: $x_1^2+x_2^2 = (9)^2 - 2(16) = 81 - 32 = 49$.
		
		Vậy, giá trị của biểu thức T là:
		$T = \dfrac{4\sqrt{17}}{49}$.
	}
\end{ex}

\begin{ex}
	Giải hệ phương trình $\begin{cases} 3x-y=5 \\ x+2y=4 \end{cases}$.
	\loigiai{
		Ta có hệ phương trình:
		$\begin{cases} 3x-y=5 \quad (1) \\ x+2y=4 \quad (2) \end{cases}$
		Từ phương trình (1), suy ra $y = 3x-5$.
		Thế $y = 3x-5$ vào phương trình (2):
		\begin{align*} x + 2(3x-5) &= 4 \\ x + 6x - 10 &= 4 \\ 7x &= 4+10 \\ 7x &= 14 \\ x &= \dfrac{14}{7} \\ x &= 2 \end{align*}
		Thay $x=2$ vào $y = 3x-5$:
		$y = 3(2) - 5 = 6 - 5 = 1$.
		Vậy, nghiệm của hệ phương trình là $(x;y) = (2;1)$.
		
		\textbf{Cách khác (phương pháp cộng đại số):}
		Nhân phương trình (1) với 2:
		$\begin{cases} 6x-2y=10 \quad (1') \\ x+2y=4 \quad (2) \end{cases}$
		Cộng vế theo vế của (1') và (2):
		$(6x-2y) + (x+2y) = 10+4$
		$7x = 14$
		$x = 2$.
		Thay $x=2$ vào phương trình (2):
		$2+2y = 4$
		$2y = 2$
		$y = 1$.
		Vậy, nghiệm của hệ phương trình là $(x;y) = (2;1)$.
	}
\end{ex}
\begin{ex}
	Một toà nhà chung cư có 30 tầng, được đánh số lần lượt từ 1 đến 30. Bạn Bình vào thang máy ở tầng 1, bấm chọn ngẫu nhiên số một tầng để đi lên. Tính xác suất của các biến cố: A: “Bình đi lên tầng có số là một số nguyên tố”.
	\loigiai{
		Bạn Bình vào thang máy ở tầng 1 và bấm chọn ngẫu nhiên một tầng để đi lên. Các tầng có thể được chọn là từ tầng 2 đến tầng 30 (vì bạn Bình đang ở tầng 1 và đi lên).
		Không gian mẫu $\Omega$ là tập hợp các tầng mà Bình có thể chọn để đi lên:
		$\Omega = \{2, 3, 4, 5, ..., 29, 30\}$.
		Số phần tử của không gian mẫu là $n(\Omega) = 30 - 2 + 1 = 29$.
		
		Biến cố A: “Bình đi lên tầng có số là một số nguyên tố”.
		Ta liệt kê các số nguyên tố trong khoảng từ 2 đến 30 (các tầng Bình có thể đến):
		Các số nguyên tố là: $2, 3, 5, 7, 11, 13, 17, 19, 23, 29$.
		Tập hợp các kết quả thuận lợi cho biến cố A là:
		$A = \{2, 3, 5, 7, 11, 13, 17, 19, 23, 29\}$.
		Số phần tử của biến cố A là $n(A) = 10$.
		
		Xác suất của biến cố A là:
		$P(A) = \dfrac{n(A)}{n(\Omega)} = \dfrac{10}{29}$.
		
		Vậy, xác suất để Bình đi lên tầng có số là một số nguyên tố là $\dfrac{10}{29}$.
	}
\end{ex}

\setcounterref{ex}{0}
